\documentclass[11pt]{article}
\usepackage[T1]{fontenc}
\usepackage[utf8]{inputenc}
\usepackage{lmodern,microtype}
\usepackage[letterpaper,margin=0.87in]{geometry}
\usepackage{amsmath,amssymb,amsthm,mathtools}
\usepackage{booktabs,longtable,array,multicol}
\usepackage{xcolor,fancyhdr,enumitem}
\usepackage[colorlinks=true,linkcolor=blue!45!black,citecolor=blue!45!black,urlcolor=blue!45!black]{hyperref}
\hypersetup{pdftitle={Lower degree and fewer witnesses for Diophantine consistency},pdfauthor={AI-assisted research prepared for Vladimir Reshetnikov and ProveIt}}
\definecolor{ink}{RGB}{25,46,68}
\pagestyle{fancy}\fancyhf{}
\fancyhead[L]{\small Diophantine consistency: trace and radical reductions}
\fancyhead[R]{\small ProveIt}
\fancyfoot[C]{\thepage}
\setlength{\headheight}{14pt}
\setlength{\parindent}{0pt}\setlength{\parskip}{0.35em}
\setlist{nosep,leftmargin=*}
\emergencystretch=2em
\allowdisplaybreaks
\newtheorem{theorem}{Theorem}[section]
\newtheorem{lemma}[theorem]{Lemma}
\newtheorem{proposition}[theorem]{Proposition}
\newtheorem{corollary}[theorem]{Corollary}
\theoremstyle{definition}
\newtheorem{definition}[theorem]{Definition}
\newtheorem{remark}[theorem]{Remark}
\newcommand{\Z}{\mathbb Z}\newcommand{\N}{\mathbb N}\newcommand{\Q}{\mathbb Q}
\newcommand{\Con}{\operatorname{Con}}\newcommand{\lcm}{\operatorname{lcm}}

\title{\color{ink}\bfseries Lower degree and fewer witnesses\\for Diophantine consistency\par\large An explicit strong-trace reduction, integral radical projections,\\and the diagonal power $t^t$}
\author{AI-assisted mathematical research\\\normalsize Prepared for Vladimir Reshetnikov and the ProveIt project}
\date{25 October 2026}
\begin{document}
\maketitle\thispagestyle{empty}
\begin{abstract}
We give two complementary improvements to fixed Diophantine equations whose
unsolvability is provably equivalent in ZFC to $\Con(\mathrm{ZFC})$.
The compact construction has ten strictly positive unknowns, exact degree
$232$, and an $87$-operation arithmetic circuit; its matched predecessor
has degree $304$ in $87$ operations. Four further compact tradeoffs are
supplied. The key observation is that the doubled strong Pell trace is
already a polynomial congruent to one modulo the main ordinate. An
auxiliary expression can therefore use this trace once instead of twice.
The proof treats negative intermediate ranks before invoking their
intended interpretations.

For unrestricted integer unknowns, we give completely expanded quartics
in $234$, $235$, and $236$ variables. The sparse $236$-variable endpoint
has $172$ quadratic residuals, $4{,}992$ collected monomials, and coefficient
height $10{,}722$. A common radical of the quadratic parts exposes
integral directions that can be eliminated without increasing degree;
six exact projections are certified in each direct endpoint. A further
semantic simplification proves that $F=t^t$ suffices for proof-label
recovery in the factorial-free binomial compiler.

All resulting equations are ordinary polynomials with fixed integer
coefficients. Complete literal data, exact transformations, executable
checks, retained arithmetic proofs, and source calibrations accompany the
paper. These are ordinary mathematical proofs composed with identified
repository results; no new end-to-end Lean/Rocq or encoded object-ZFC proof
is claimed. The comparisons concern an immutable inspected snapshot,
and do not establish a globally simplest equation or literature-wide priority.
\end{abstract}
\vfill
\begin{center}\small
Source snapshot: \texttt{b7e35224c04fe1031c0a53f80a395e9b32205022}.\\
Positive and unrestricted integer domains are kept distinct.
\end{center}
\newpage\tableofcontents\newpage

\section{The result and the comparison standard}

A small number of unknowns can make a formula short while requiring
enormous fixed coefficients and high degree. A direct encoding of finite
proof certificates can keep all coefficients small and the degree four,
at the cost of more unknowns. We improve both directions without collapsing
their different measures into a single unsupported minimum claim.

Our two main displayed equations are
\begin{align}
 P_{\mathrm{tr}}(J,F,\alpha,i,T,w,s,n,\eta,\zeta)&=0,
 &\boldsymbol u&\in\Z_{>0}^{10},\label{eq:compact-statement}\\
 P_{236}(x_0,\ldots,x_{235})
  =\sum_{j=0}^{171}R_j(\boldsymbol x)^2&=0,
 &\boldsymbol x&\in\Z^{236}.\label{eq:direct-statement}
\end{align}
Section~\ref{sec:compact} defines every polynomial abbreviation in the first
equation. Appendix~\ref{app:direct} prints every quadratic $R_j$ in the
second. The archive also contains complete expansions of all three direct
quartics, with ordered coordinates and integer coefficients. The four fixed
compact numerals are specified by the retained finite recipe and exact
binary payloads; they are not adjustable parameters.

\begin{theorem}[Consistency equivalences]\label{thm:calibration}
For the literal equations supplied in this release, the new proofs below,
composed with the identified inherited arithmetic and syntactic calibration,
establish
\begin{align*}
 \mathrm{ZFC}&\vdash
 [\neg\exists\boldsymbol u\in\Z_{>0}^{10}\ P_{\mathrm{tr}}(\boldsymbol u)=0]
 \longleftrightarrow\Con(\mathrm{ZFC}),\\
 \mathrm{ZFC}&\vdash
 [\neg\exists\boldsymbol x\in\Z^{d}\ P_d(\boldsymbol x)=0]
 \longleftrightarrow\Con(\mathrm{ZFC}),\quad d\in\{234,235,236\}.
\end{align*}
The same positive-domain statement holds for the other four supplied
compact circuits. No consistency assumption is used in deriving these
equivalences.
\end{theorem}

This is an equivalence provable in ZFC, not a proof in ZFC of either side.
Each equation has a solution precisely when a finite refutation exists,
under its stated encoding. Finite experiments with candidate solutions
would not establish consistency.

\subsection{Measured changes}
\begin{table}[ht]\centering\small
\begin{tabular}{@{}rrrr@{}}\toprule
Unknowns/domain & Operations & Parent degree & New degree\\\midrule
10 positive & 75 & 626 & 402\\
10 positive & 76 & 464 & 352\\
10 positive & 79 & 382 & 290\\
10 positive & 80 & 338 & 258\\
10 positive & 87 & 304 & 232\\
\bottomrule\end{tabular}
\caption{The five compact improvements. Each retains its parent's operation count and complete fixed coefficient tuple.}\label{tab:compact}
\end{table}
\begin{table}[ht]\centering\small
\begin{tabular}{@{}lrrrr@{}}\toprule
Integer quartic & Variables & Rows & Terms & Height\\\midrule
Sparse diagonal power & 236 & 172 & 4,992 & 10,722\\
Dense diagonal power & 235 & 171 & 10,536 & 10,730\\
Shared factorial error & 234 & 171 & 21,385 & 44,800\\
\bottomrule\end{tabular}
\caption{Three new direct equations over unrestricted integers. All have exact degree four.}\label{tab:direct}
\end{table}
The matched direct parents have respectively $242$, $241$, and $240$ variables.
The first two new branches also replace the label base by $t^t$.
Each final radical projection phase removes exactly six variables and six rows.
The $240$-variable shared-error parent has $21{,}482$ terms and height $6{,}144$;
its $234$-variable descendant reduces dimension and support but increases height.


Every operation count charges a live binary addition, subtraction, or
multiplication, including a square and multiplication by a fixed scalar.
Repeatedly used named values are shared. Such names are not new existential
unknowns. Degree means total degree after all abbreviations are substituted,
and the finite verifier certifies an attained highest homogeneous form,
not merely a syntactic upper bound. A direct polynomial's term count
refers to its fully collected sum of squares; its coefficient height is
the largest absolute coefficient in that expansion.

The compact results keep all four old numerical ports intact. Their
degree and evaluation improvements therefore do not hide a coefficient
change. The direct results expose their support/height tradeoffs. The
$234$-variable endpoint saves the most coordinates among the three direct
forms, while $236$ is much sparser. Neither comparison changes the domain.

\section{The compact equation}\label{sec:compact}

All ten independent variables in \eqref{eq:compact-statement} are strictly
positive. Let $\beta,K_0,M_C,M_F$ be the complete inherited BFS943 compact44
tuple, with its corrected initializer and source convention retained. Its
fixed numerical bounds include
\[
 \beta\ge31\text{ odd},\qquad K_0>0,\qquad
 0<M_C<\beta<M_F<2\beta.
\]
These bounds help prove the arithmetic reduction, but do not alone select
the ZFC instance. The full fixed tuple in the archive selects that instance.
Define
\begin{gather*}
 h=\beta J,\qquad q=h+1,\qquad C=q-F-\alpha,\\
 X=wq,\qquad Y=qs,\qquad k=\eta+\zeta,\qquad c=kY+\eta,\\
 A=Y(X+1)+1,\qquad \Delta=A^2-1,\qquad H=4A-5,\\
 v=XY^2,\qquad E=1+v(v+1)k^2,\\
 R=[h(q-F)+\alpha](q^2-1)+(M_C+qM_F)J,\\
 B=(K_0+w)C-F,\qquad U=ic^2,\\
 a=\Delta U^2E,\qquad G=E(1+a),\qquad f=1+2a,\\
 Q=4\Delta^2U^2G,\qquad b=cT-R,\qquad V=fb-c,\\
 L=bV,\qquad K=X(3c-2X)+2,\\
 e_6=6(k-R-1)-v(R+1)(R-1)(R+3),\\
 D_6=6q^6,\qquad M=hD_6H,\\
 Z=hD_6[H(n-LB)+K]-H[D_6B+h(e_6-6q^4s)].
\end{gather*}
The promised polynomial is
\begin{equation}\label{eq:compact}
 \boxed{\quad P_{\mathrm{tr}}
 =Q\bigl(Z^2-M^2GV^2\bigr)+M^2G-Z^2.\quad}
\end{equation}
All these expressions are ordinary polynomial abbreviations. The
verified circuit evaluates exactly \eqref{eq:compact}, exploiting
$Q=4\Delta a(1+a)$ and the shared products listed in its JSON definition.
Its exact count is $87$, and its exact total degree is $232$.

The lower-case $f$ is the explicit strong Pell trace. It differs from
the independent uppercase variable $F$ in the coefficient mask. Neither
should be identified with the $F=t^t$ variable in the separate direct
construction.

\subsection{The algebraic opportunity}

The matched predecessor sets $V=f^2b-c$, evaluating $f^2$ through the
identity $1+4a(1+a)$. We instead set
\begin{equation}\label{eq:tracechange}
                  V=f(cT-R)-c.
\end{equation}
Since $U=ic^2$, we have $c\mid a$ and therefore
\begin{equation}\label{eq:residues}
 f\equiv1\pmod c,\qquad V\equiv-c\pmod f,\qquad V\equiv-R\pmod c.
\end{equation}
The old proof needs precisely the latter two residues at its index-recovery
stage. The special doubled-trace formula supplies $f\equiv1\pmod c$,
whereas the previous expression only exploited $f^2\equiv1\pmod c$.

A congruence observation alone would be insufficient. The proof also
needs absolute size estimates and positive witness construction. Some
intermediates can be negative at arbitrary positive inputs. The following
sections verify these obligations for the changed expression.

\section{Soundness of the trace reduction}

We retain the proved Pell folding and native decoding interfaces from the
identified source packets, and state where they enter. Their supporting
proofs and calibrations are included in the release. We use
\[
 \chi_A(m)+\psi_A(m)\sqrt{A^2-1}
       =(A+\sqrt{A^2-1})^m\qquad(m\ge0).
\]

\subsection{Recover roots before interpreting the mask}

Positive inputs imply $q\ge32$, $X,Y\ge q$, $k\ge2$, $c>kY$,
$A\ge3$, $U\ge c^2$, $a>0$, $f>2c$, $Q>1$, and $M>0$.
For integral $b$, the value $V=fb-c$ is nonzero: if $b=0$ then
$V=-c$; otherwise $|V|\ge f-c>1$.
At a zero of \eqref{eq:compact}, write $z=Z/M$ to obtain
\begin{equation}\label{eq:productnorm}
                    (Q-1)z^2=G(QV^2-1).
\end{equation}
Since $Q=4\Delta^2U^2G$, both $Q-1$ and $QV^2-1$ are coprime to
$G$. Prime valuations force each exponent in $G$ to be even. Thus
$G=g^2$ with $g>0$. Moreover $\gcd(E,1+a)=1$, because $E\mid a$,
so $E=\tau^2$, $1+a=f_0^2$, and $g=\tau f_0$ for positive integers.

The inherited rational adjacent-Pell descent, applied to
$QV^2-(Q-1)y^2=1$ with $y=|z|/g$, proves $y\in\Z_{>0}$.
This is an integrality theorem, not an assumption that a rational square
root is integral. Thus $z=\varepsilon gy\in\Z$ for
$\varepsilon\in\{-1,1\}$. Put $I=2Ug$. The exact identities
\begin{equation}\label{eq:strong}
 f^2-\Delta I^2=1,\qquad Q=(\Delta I)^2,\qquad c^2\mid I
\end{equation}
are unchanged. The signed folding interface uses \eqref{eq:strong},
the adjacent norm, $c\le I$, and $V\equiv-c\pmod f$; it does not use
the extra trace factor in the old $V$. It yields
$c=\psi_A(p)$ for a positive odd $p$.
The first norm gives $k=2\psi_{2v+1}(m)$ for $m\ge1$.
Index $m=1$ would force $k=2$, $\eta=\zeta=1$, and
$c=2Y+1$, strictly between the first two positive main ordinates.
Hence $m\ge2$ and
\begin{equation}\label{eq:earlyc}
              c>kY\ge(8v+4)Y>8XY^3\ge8q^4.
\end{equation}
In particular $c>X$ and $K>0$. No positivity of $R,C,B,b$, or $V$
has been assumed.

\subsection{The absolute estimate needed for transport}

If $b=0$, the adjacent norm would give
$y^2-c^2=(c^2-1)/(Q-1)$ strictly between zero and one.
Here $Q=(\Delta I)^2\ge9c^4>c^2$ by \eqref{eq:strong}.
Thus $b\ne0$, and $b,V$ have the same sign. With $D_0=\Delta U^2$,
\[
 \frac{a^2}{g^2}=\frac{D_0a}{1+a}>\frac{D_0}{2}>9.
\]
Consequently $a>3g$, $|V|\ge f-c>a>3g$, $|V|>|b|$, and
$y<3|V|/2$.

The exclusion of main index $p=3$ is independent of the formula for
$V$. If $c=4A^2-1$ and $m\ge3$, Pell growth gives $k>32v^2$,
incompatible with $c=kY+\eta$. Thus $m=2$ and $k=8XY^2+4$.
For $u=X+2-2Y$, expansion gives
\[
 \eta=4Y^2(Xu+1)+8Y(X+1)+3-4Y.
\]
If $u\le-1$ this is negative; if $u\ge2$ it exceeds $k$.
Both contradict $0<\eta<k$. The remaining cases force
$2Y-X\in\{1,2\}$, impossible since $q\mid X,Y$ and $q\ge32$.
Hence $c>c_3:=4\Delta+3$.

Classification of the adjacent norm, with $S=\Delta I$, gives
$|V|=\chi_S(j)/S$ for odd $j\ge1$. Index one is excluded.
Index three would give $|V|=4Q-3$. As $Q=\Delta(f^2-1)$,
either sign would imply one of
\[
 f(b-4\Delta f)=c-c_3,\qquad f(b+4\Delta f)=c+c_3.
\]
The right sides are strictly between zero and $2c<f$.
Thus $j\ge5$ and
\[
 |V|\ge16Q^2-20Q+5>8Q^2,\qquad
 |b|>\frac{8Q^2-c}{f}>8f^3-\frac12>4g.
\]
Here $Q>f^2$ follows from $\Delta\ge3$ and $f\ge2$.
Therefore $L=bV=|b||V|$ satisfies
\begin{equation}\label{eq:transportbound}
 L>0,\qquad gy<\frac38L<\frac12L,\qquad
 \frac{|b|}{q^6}<\frac12L.
\end{equation}

\subsection{Negative ranks in the Lucas-corrected equation}

For proof only, put
\[
 t_0=B/h,\quad W=1+hL,\quad\kappa=K/H,\quad
 \sigma=s/q^2,\quad\bar e=e_6/6,\quad\rho=\bar e/q^6.
\]
The carrier identity gives
\begin{equation}\label{eq:transport}
 n=\varepsilon gy-\kappa-\sigma+Wt_0+\rho.
\end{equation}
Initially $\bar e$ may be rational. Write $B_R=(R+1)(R-1)(R+3)$.
Since $cT>k-1$,
\[
 \rho<\frac b{q^6}-\frac{vB_R}{6q^6}.
\]
If $B_R\ge0$ and $B\le-1$, then $Wt_0<-L$; positivity of
$\kappa,\sigma$ and \eqref{eq:transportbound} force $n<0$.

The adverse sign $B_R<0$ occurs precisely at integer ranks $R\le-4$
and $R=0$. For $R=-r\le-4$, the exact mask identity
\[
 R=(hC+q\alpha)(q^2-1)+(M_C+qM_F)J
\]
gives $C<0$, $|C|>r/q^3$, and $B\le-1-w|C|$.
The adverse correction is less than $ws^2r^3/(6q^3)$.
But $b=cT+r>r,c$, so $V>(f-1)b$ and
$L>(f-1)r^2$, while $f-1=2a>s^2$.
The additional term $-w|C|L$ therefore dominates this correction.
The remaining terms again yield $n<0$ by \eqref{eq:transportbound}.
For $R=0$, we have $|C|\ge1$ and $L>f-1>s^2$; the adverse
term $ws^2/(2q^3)$ is likewise dominated. This covers every integer rank.

It follows that $B\ge0$ and $C>0$, because the independent $F$ is
positive. The inherited mask estimate is now licensed:
$3q<R<q^4-2$. Equation~\eqref{eq:earlyc} gives $c>2R$, so
$b,V>0$. The residues \eqref{eq:residues} and the strong pair
\eqref{eq:strong} now meet all hypotheses of the normalized rank and
signed auxiliary-index lemmas. They yield $p=R$ and the old parity
conclusions.
The other size premise $c>2p$ follows from $c=\psi_A(p)$,
$A\ge3$, and the already established $p\ge5$.

\subsection{The unchanged decoding stages}

The identity
\[
 2\psi_{1+2v}(m)=2m+8v\binom{m+1}{3}+32v^2Z_m,\qquad Z_m\in\N,
\]
and the carrier recover the first index and remaining denominators.
Their data $q,X,Y,k,c,R$ and integral carrier $z$ have not changed.
The inherited proof first obtains divisibility by $q^4$, bounds the
index difference by $R<q^4$, and only then obtains divisibility by
$q^6$. Its valuation calculation covers the primes two and three
dividing the factor six. The whole-product projection, native witnesses,
and fixed-source refutation calibration then apply. The retained source
proof supplies these unchanged stages; no new cancellation modulo a
possibly nonunit factor is introduced.

\section{A polynomial completeness map}

There is a concrete map from every old positive zero, stronger than
rebuilding canonical witnesses alone. Define the polynomial
\[
 d=\frac{f-1}{c}=2\Delta i^2c^3E.
\]
Given an old positive zero with $V_{\rm old}=f^2b_{\rm old}-c$, set
\begin{equation}\label{eq:Tnew}
 T_{\rm new}=fT_{\rm old}-dR.
\end{equation}
All other independent inputs except $n$ are fixed. Then identically
\[
 b_{\rm new}=f b_{\rm old},\qquad V_{\rm new}=V_{\rm old}.
\]
For the $L=bV$ charts, set
\[
 n_{\rm new}=n_{\rm old}+(f-1)L_{\rm old}B.
\]
For the $L=V^2$ chart, leave $n$ unchanged. In either case
$n_{\rm new}-L_{\rm new}B=n_{\rm old}-L_{\rm old}B$, so the entire
carrier and final polynomial are preserved under this substitution.

At an old positive zero, the old proof gives $0<R<c/2$ and $B\ge0$.
Since $T_{\rm old}\ge1$,
$T_{\rm new}>f-(f-1)/2>0$; it is integral because $d$ is a polynomial.
The old $b,V,L$ are positive, so $n_{\rm new}>0$ as well.
This proves completeness on the required domain. Soundness for arbitrary
new positive zeros was proved separately above. We do not claim an
integral inverse of this map on every new zero.

The same trace replacement works in four simpler supplied charts.
Their complete formulas and inherited source proofs are retained in
the compact component; the altered transport estimates are no stronger
than those just established. Table~\ref{tab:compact} records all five
degree/cost tradeoffs.

\section{A smaller arithmetic certificate base: \texorpdfstring{$F=t^t$}{the diagonal power}}\label{sec:selfpower}

This section belongs to the direct proof-calculus construction and uses a
different $F$ from the compact mask. The factorial-free predecessor uses
$F=[2(t+1)]^t$ to separate proof-label multiplicities and make its binomial
CRT modulus cover every row. The diagonal power $F=t^t$ suffices for both
jobs. Its exact power relation is compiled into polynomial Pell residuals;
it is not left in the final Diophantine equation.

\begin{lemma}[Label balance at the diagonal power]\label{lem:labels}
Let $t\ge1$, $F=t^t$, and $W_j=1+jF$ for $0\le j\le t$.
Suppose $v_i,p_i,q_i\in\{0,\ldots,t\}$ for $0\le i<t$, and
\begin{equation}\label{eq:labels}
 \prod_{i<t}W_{v_i}=(1+F)\prod_{i<t}W_{p_i}W_{q_i}.
\end{equation}
At every positive label $j$, its conclusion multiplicity minus its
premise multiplicity is one at $j=1$ and zero otherwise.
\end{lemma}
\begin{proof}
Discard the label-zero factors, which are one. If the two factor counts
differ, their smaller count $r$ is at most $t$, because the left side
has at most $t$ nonunit factors. The side with more factors is at least
$(F+1)^{r+1}$, and the other side is at most
$(1+tF)^r\le t^r(F+1)^r$. But $F+1=t^t+1>t^r$, a contradiction.
Hence their counts agree, with common value $r\le t$.

Replace $F$ by an indeterminate $z$. Each side is a product of $r$
linear factors $1+jz$, where $1\le j\le t$. Its coefficient of $z^k$,
for $0\le k\le r$, is positive and at most
\[
 \binom rk t^k\le t^{r-k}t^k=t^r\le F.
\]
For the inequality, inject the $(r-k)$-element subsets of an $r$-element
set into ordered $(r-k)$-tuples from a $t$-element set. Thus every
coefficient of the difference polynomial has absolute value strictly
less than $F$. If the difference were nonzero, divide its evaluation
at $F$ by the power corresponding to its lowest nonzero coefficient
and reduce modulo $F$. That coefficient would be a nonzero multiple
of $F$ of absolute value less than $F$. This is impossible.
The products are therefore equal in $\Z[z]$, and the distinct roots
$-1/j$ give the claimed multiplicities. When $F=1$, the coefficient
difference bound itself already implies that the difference is zero.
\end{proof}

The strict bound is on the \emph{difference} coefficients. Bounding
each entire product by its coefficient sum misses this improvement
and retains the unnecessarily larger base $(t+1)^t$.

\begin{lemma}[Single-conclusion flow]
Take a finite multiset of valid Horn rules with one conclusion and at most
two premises; pad unused premises with freely derivable truth, label zero.
If their multiplicities have the balance in Lemma~\ref{lem:labels},
then label one has a finite derivation.
\end{lemma}
\begin{proof}
Form the closure of truth under these finitely many rule instances.
If label one were outside, every rule whose conclusion is outside would
have at least one premise outside. The sum of conclusions minus premises
over outside labels would be at most zero. The prescribed balance
instead gives one, a contradiction.
\end{proof}

This inherited flow argument explains why an unordered certificate with
cycles does not create a false proof. It needs no ordering of the rows.

\subsection{The binomial compiler still covers every row}

Let $\Lambda_t=\lcm(1,\ldots,t)$ and $\mathcal R=F^2$. We need only
\emph{numerical domination}, not $\Lambda_t\mid F$:
\[
 \Lambda_t\le t!\le t^t=F,\qquad \mathcal R\ge\Lambda_t^2.
\]
Indeed $\Lambda_t$ need not divide $t^t$: three divides $\Lambda_{16}$
but not $16^{16}$. Avoiding this divisibility requirement was the
purpose of the inherited factorial-free compiler.

For clarity, its two prime-power facts are recalled here. Put
$\mathcal B=\binom{\mathcal R+t}{t}$. Every full prime power
$p^e\Vert\mathcal B$ divides some $\mathcal R+j$, $1\le j\le t$:
Legendre's valuation formula consists of zero-one summands and its
sum is at most $\max_jv_p(\mathcal R+j)$. Conversely, each row $j$
has a full prime power of $\mathcal B$ exceeding $t$. Since
$\mathcal R+j>\Lambda_t^2$, some $v=v_p(\mathcal R+j)$ exceeds
$2\lfloor\log_p t\rfloor$. For this prime, the valuation in
$\mathcal B$ is exactly
\[
 v-v_p\!\left(t\binom{t-1}{j-1}\right).
\]
The subtracted valuation is at most $\lfloor\log_p t\rfloor$, because
$t\binom{t-1}{j-1}\mid\Lambda_t$. This classical denominator bound
also follows from Farhi's binomial lcm identity~\cite{farhi}.
The resulting full prime power therefore exceeds $t$ and identifies
the row uniquely. The first fact covers \emph{every prime} in the
converse CRT construction, while the second covers \emph{every row}
in soundness. Both remain valid with $F=t^t$.

\subsection{Parity and cofinal completeness}

The retained horizon is $t=16+h_0^2$. For completeness choose the
cofinal subfamily with even $h_0$, so $t$ is even. Then
\[
 v_2(F)=t\,v_2(t)\ge t.
\]
Thus $F$ and the beta base $y=F(2u+1)^2$ are even.
Let $O$ be the odd part of $t!$. Choose $2u+1$ to be a sufficiently
large positive odd multiple of $O$. Since $v_2(F)\ge t>v_2(t!)$,
this makes $y$ an even multiple of $t!$, as required by the inherited
pairwise-coprime beta-modulus construction. This choice supplies the
missing odd prime factors; it does not incorrectly assume $t!\mid t^t$.
The beta moduli $1+(j+1)y$ are all odd, so their CRT progression
admits arbitrarily large even representatives for the beta code.
The binomial modulus $\mathcal B$ is odd: since
$v_2(F^2)>\lfloor\log_2t\rfloor$,
$v_2(F^2+j)=v_2(j)$ for $1\le j\le t$, and the numerator and $t!$
have equal two-adic valuations. All inherited CRT parity choices
remain available.

Soundness imposes no new evenness condition on an arbitrary satisfying
horizon. The independent audit examines each use of the old base
$[2(t+1)]^t$, rather than replacing it in the formula and assuming
all estimates survive. The needed size bound follows already from
$F=t^t\ge16^{16}=2^{64}$ at $t\ge16$, together with
$F\ge\Lambda_t$. The source edit also changes the actual power-call
base and the shared frame metadata, which are checked separately.

\section{Eliminating common quadratic directions}\label{sec:radical}

The second direct improvement is algebraic and reusable. An ordinary
monic substitution can raise degree: replacing a variable by a quadratic
expression where it appears quadratically elsewhere may create quartic
residuals and an eighth-degree sum of squares. Instead, we find directions
along which \emph{every} residual's quadratic part is constant.

Write
\[
 R_j(\boldsymbol x)=c_j+\ell_j(\boldsymbol x)+Q_j(\boldsymbol x),
\]
where $\ell_j$ and $Q_j$ are homogeneous of degrees one and two.
Let $H_j$ be the integer Hessian matrix of $Q_j$; a term $a x_r^2$
contributes $2a$ in its diagonal position.

\begin{definition}
The common quadratic radical is
\[
 \mathcal K=\bigcap_j\ker_{\Q}H_j\subseteq\Q^n.
\]
\end{definition}

\begin{theorem}[Integral radical projection]\label{thm:radical}
Suppose $\boldsymbol d\in\Z^n\cap\mathcal K$ has $d_v=1$ at a
coordinate $v$, and $\ell_r(\boldsymbol d)=\varepsilon\in\{1,-1\}$
for a residual $r$. The common integer zero set of the system is then
in explicit polynomial bijection with the common integer zero set
of $m-1$ quadratic residuals in $n-1$ coordinates.
\end{theorem}
\begin{proof}
Make the integral coordinate change
\[
 x_v=y_v,\qquad x_w=y_w+d_wy_v\quad(w\ne v).
\]
Its inverse is $y_v=x_v$, $y_w=x_w-d_wx_v$, so it is unimodular.
Since $H_j\boldsymbol d=0$, the quadratic part of every transformed
residual is independent of $y_v$. Each transformed residual is
\[
 S_j(\widehat{\boldsymbol y})+\ell_j(\boldsymbol d)y_v,
\]
where the hat omits coordinate $v$ and $S_j$ is quadratic.
The selected row has the unique integral solution
$y_v=-\varepsilon S_r$. Substituting in all other rows gives
\[
 S_j-\varepsilon\ell_j(\boldsymbol d)S_r,
\]
still quadratic. Delete the selected row and coordinate.
Restriction maps an old zero to a new zero. Restoring $y_v$ by
the displayed expression and undoing the shear gives its unique
inverse. All maps have integer coefficients.
\end{proof}

The sum of squares still has degree at most four. Its polynomial need not equal the
old polynomial; the zero sets correspond. This distinction matters
when comparing collected term counts and coefficient heights.

\subsection{A small example}

For
\[
 R_1=(x+y)^2+z-1,\qquad R_2=(x+y)^2+x+2y,
\]
the direction $(1,-1,0)$ belongs to both Hessian kernels. Substituting
$x=u$, $y=v-u$ gives $R_1=v^2+z-1$ and $R_2=v^2+2v-u$.
The second row restores $u=v^2+2v$, leaving only $v^2+z-1=0$.
Both the shear and the restoration are integral. The large certificates
use exactly this mechanism.

\subsection{The six actual projections}

At every stage, the supplied certificate records the ordered variable
names, the integer radical direction, the unit coordinate, the selected
row, the shears, the quadratic restoration, and the retained-coordinate
map. Each branch uses six such projections. Each step deletes exactly
its selected coordinate and row; no further unused variable is discarded
at this stage.

The independent checker verifies $H_j\boldsymbol d=0$ for every row,
the unit slope, both coordinate identities, the restored pivot equation,
and every substituted residual as an exact polynomial identity.
It expands and collects the final sum of squares. Earlier preprocessing
may remove unused coordinates; that separate existence-of-zeros step
is not described as a bijection of all raw assignments.

The search computes rational kernels and clears denominators, but its
output certificate is integral. No search heuristic is a mathematical
premise. Exact modular rank certificates show that the final stacked
Hessian matrix has full column rank modulo $101$ in all three branches.
The corresponding integer minor is nonzero, so the common radical is
zero over $\Q$. There are no further directions of this particular type
at these charts. This is not a global minimality statement.

\section{From proof calculus to the fixed integer quartics}

The direct branch uses the numerical proof calculus $C_{222}$. The
retained source compiles its finite numerical Horn rule families, an
unordered flow certificate, bounded range assertions, binomial CRT
row selection, and a synchronized Pell power representation. The
result is a system of quadratic residuals with fixed integer coefficients.
It has no Turing machine coefficient tuple and does not invoke Jones's
universal polynomial.

The $235$- and $236$-variable branches begin with the factorial-free
compiler~\cite{factorialfree}, change the power request to $t^t$ by
Section~\ref{sec:selfpower}, and apply the six radical projections.
Their support tradeoff reflects an earlier dense affine elimination
used only in the $235$ branch. The $234$-variable branch applies the
six projections to the shared-factorial-error quartic~\cite{sharederror};
it does not claim to remove that branch's factorial identity.

The relevant chain, formalizable internally in ZFC, is
\[
 \begin{gathered}
  \text{an integer zero of the final quartic}\\
  \Longleftrightarrow\text{a valid compiled arithmetic certificate}\\
  \Longleftrightarrow\text{a finite }C_{222}\text{ refutation}\\
  \Longleftrightarrow\neg\Con(\mathrm{ZF})
  \Longleftrightarrow\neg\Con(\mathrm{ZFC}).
 \end{gathered}
\]
The last step is syntactic equiconsistency, using the retained
constructible-model calibration. It does not assume a transitive
model or an externally standard model. The archived calibration
review identifies the source interfaces on both sides.

An external proof-assistant theorem about standard numerical objects
and an encoded object-language ZFC derivation are different artifacts.
This release supplies ordinary arguments formalizable in ZFC and
retains the inherited formal-interface sources. It does not claim
that every stage has been newly elaborated as a Lean, Rocq, or
object-ZFC proof term. A finite Python replay authenticates coefficients
and transformations; it does not replace unbounded Pell, CRT, and
syntactic calibration arguments.

\section{Verification and reproducibility}\label{sec:verification}
The final release checks are organized by mathematical obligation.
\begin{itemize}
\item The compact checker authenticates $105$ inherited source files and
all four fixed numerical payloads. For each of the five circuits it checks
topology, liveness, the complete parent/child gate relation, and every
polynomial stage through $32$--$37$ exact cuts. Nine auxiliary identities
include the strong-trace algebra and polynomial completeness map.
\item It computes the complete highest homogeneous forms. Eighteen
deliberate operator corruptions per circuit are rejected, for a total of
$90$. Normal and optimized-Python receipts agree. Exact reconstruction
of the fixed tuple and byte-for-byte circuit rebuilding pass.
\item The direct checker independently represents monomials as
variable/exponent pairs. It checks the actual diagonal-power source edit,
$18$ Horn identities, $202$ predecessor monic projections, the selected
affine maps, and all six new radical projections for each endpoint.
\item It independently expands every final quartic, verifies every
coefficient, checks all reported metrics, and verifies the terminal
Hessian ranks modulo $101$. Five deliberate corruptions of directions,
restorations, coordinate deletion, residuals, and expanded coefficients
are rejected.
\item The scoped predecessor checks pass both normally and with Python
assertions disabled. A disposable reconstruction reproduces $28$
selected mathematical JSON files byte for byte. These retained
dependency subsets omit historical PDFs and obsolete experiments;
their scoped entry points are explicitly documented.
\item A root release replay checks the component verifiers and ensures
that the report's literal residual appendix, comparison tables, and
degree examples match their canonical data. The PDF is rendered and
visually inspected before packaging.
\end{itemize}

The release uses Python 3.10 or newer and its standard library for
these exact checks. Run from the extracted package directory:
\begin{verbatim}
python3 -B verify.py
python3 -B verify.py --rebuild
sh build_report.sh
\end{verbatim}
The first command checks the frozen artifacts without changing them.
The optional second additionally regenerates the selected artifacts
in disposable locations. The last rebuilds the paper from its TeX
source. Internal independent mathematical reviews are retained in
\texttt{review/}; they are AI-assisted reviews, not external human
peer review.


The canonical equations use ordered sparse terms:
$[a,[i_1,\ldots,i_r]]$ denotes $a x_{i_1}\cdots x_{i_r}$.
Repeated indices are repeated factors, and an empty index list is a
constant. Every array is finite. The evaluator accepts an ordered tuple
of exact integers; it is not a solver or a consistency decision procedure.

Every compact coefficient is included with its inherited exact
reconstruction. The actual native tuple, not an arbitrary tuple
satisfying a few bounds, is bound to the consistency statement.
Checks distinguish coefficient reconstruction, literal circuit identity,
degree certification, expansion, transformation correctness, and the
unbounded mathematical proof.

The source snapshot and archive hashes are in the ledger. The final
manifest authenticates the deliverable bytes. The TeX source builds
with a standard TeX Live installation; font versions and PDF metadata
can change the PDF hash without altering the canonical equations.

\section{What this iteration adds, and further research}

The compact advance is the use of the explicit strong-trace congruence
$f\equiv1\pmod c$, together with a complete signed proof, to replace
$f^2(cT-R)-c$ by $f(cT-R)-c$. This lowers degree at the same binary operation count: the new trace
is evaluated as $f=a+(1+a)$ using the already shared $1+a$. The new polynomial parent-to-child map gives a further
algebraic check on completeness.

The direct advance is a reusable integral projection mechanism that
finds six hidden removable directions. The $t^t$ lemma separately
simplifies an actual power request and its justification. The projection
applies to more than one inherited source, and the label lemma is
proved before any quadratic projection.

The following questions are concrete directions for subsequent work.
No unproved answer is used in this release.
\begin{enumerate}
\item Can a weaker auxiliary congruence yield a further compact
trace deletion? A new index argument would be needed: the present
residue modulo $f$ does essential work.
\item Can a changed residual presentation expose more removable
directions after the common radical becomes zero? Constant invertible
row operations preserve the common radical, so they cannot suffice alone.
\item Can a structured nonlinear coordinate change preserve quadratic
residuals and expose more unit coordinates? A valid proposal needs an
integral inverse or an explicit existential restoration.
\item Can a simpler proof calculus reduce the direct front end enough
to compensate for its arithmetic overhead and calibration? Counting
rule schemata alone does not measure the final polynomial.
\item Can the synchronized power compiler use the diagonal nature
of $F=t^t$ more deeply than changing its base? Soundness must still
cover arbitrary integer assignments.
\item Can the radical-projection and trace-congruence lemmas first
be formalized as reusable components, then connected to the internal-ZFC
calibration? This would address the actual formal gap rather than
recheck only finite numerical cases.
\end{enumerate}


\appendix
\section{Complete residuals of the sparse integer equation}\label{app:direct}
Every $x_0,\ldots,x_{235}$ below ranges over $\Z$. The equation is exactly
$\sum_{j=0}^{171}R_j^2=0$. This appendix is generated from the canonical
residual JSON and checked against those literal coefficients. The other
two residual lists and all three collected quartics are in the archive.
% Generated from direct/equations/self_power236/equation.json.
\begingroup\small
\setlength{\jot}{1pt}
\begin{align*}
R_{0} &={} -x_{3}^{2} +x_{4} \\[3pt]
R_{1} &={} x_{0} -x_{0}\,x_{2} -x_{2} +x_{8} \\[3pt]
R_{2} &={} -x_{2}\,x_{6} -x_{3} +x_{7}\,x_{8} \\[3pt]
R_{3} &={} -1 +x_{1} -x_{7}\,x_{8} \\[3pt]
R_{4} &={} 1 +x_{4} -x_{5}\,x_{19} +x_{17} \\[3pt]
R_{5} &={} -1 -x_{9}^{2} -2\,x_{9}\,x_{10} -x_{10}^{2} +x_{20} \\[3pt]
R_{6} &={} -10\,x_{9} +x_{9}^{2} -6\,x_{9}\,x_{10} -8\,x_{9}\,x_{20} -20\,x_{10} -11\,x_{10}^{2} -16\,x_{10}\,x_{20} -15\,x_{20} -4\,x_{20}^{2} \\*
&\quad{} +x_{21} \\[3pt]
R_{7} &={} -2 +2\,x_{9} -x_{9}^{2} -4\,x_{9}\,x_{21} -4\,x_{21} -4\,x_{21}^{2} +x_{22} \\[3pt]
R_{8} &={} 1 +6\,x_{9} -36\,x_{9}^{2} -24\,x_{9}\,x_{22} +2\,x_{22} -4\,x_{22}^{2} +x_{23} \\[3pt]
R_{9} &={} -2 -x_{10} -x_{10}^{2} -2\,x_{10}\,x_{11} -x_{11} -x_{11}^{2} +x_{24} \\[3pt]
R_{10} &={} -1 +x_{9} -x_{9}^{2} -2\,x_{9}\,x_{11} +x_{11} -x_{11}^{2} +x_{25} \\[3pt]
R_{11} &={} -16 +14\,x_{9} -2\,x_{9}^{2} +8\,x_{9}\,x_{11} -8\,x_{9}\,x_{24} -8\,x_{11} +8\,x_{24} -4\,x_{26}^{2} +x_{27} -2\,x_{44} \\[3pt]
R_{12} &={} -2 -8\,x_{9} +8\,x_{9}\,x_{24} +8\,x_{9}\,x_{25} -8\,x_{10} +8\,x_{10}\,x_{24} +8\,x_{10}\,x_{25} -4\,x_{20} +8\,x_{24} -4\,x_{24}^{2} \\*
&\quad{} -8\,x_{24}\,x_{25} +8\,x_{25} -4\,x_{25}^{2} +x_{28} \\[3pt]
R_{13} &={} 6 +4\,x_{9} +32\,x_{9}\,x_{10} +16\,x_{9}\,x_{20} +16\,x_{9}\,x_{25} +8\,x_{9}\,x_{28} +24\,x_{10} -32\,x_{10}\,x_{25} -16\,x_{10}\,x_{28} \\*
&\quad{} +12\,x_{20} -16\,x_{20}\,x_{25} -8\,x_{20}\,x_{28} -x_{21} +8\,x_{25} -16\,x_{25}^{2} -16\,x_{25}\,x_{28} +4\,x_{28} -4\,x_{28}^{2} \\*
&\quad{} +x_{29} \\[3pt]
R_{14} &={} 7 -28\,x_{9} -8\,x_{9}\,x_{11} +16\,x_{9}\,x_{20} +8\,x_{9}\,x_{25} -27\,x_{10} -9\,x_{10}^{2} +20\,x_{10}\,x_{11} -16\,x_{10}\,x_{24} \\*
&\quad{} +12\,x_{10}\,x_{25} -32\,x_{11} -28\,x_{11}^{2} +24\,x_{11}\,x_{20} -24\,x_{11}\,x_{24} +8\,x_{20}\,x_{25} +x_{21} -32\,x_{24} \\*
&\quad{} -4\,x_{24}^{2} -x_{26} +x_{30} +x_{44} \\[3pt]
R_{15} &={} 34 -2\,x_{9} -8\,x_{9}^{2} -16\,x_{9}\,x_{10} -34\,x_{9}\,x_{11} -16\,x_{9}\,x_{20} -4\,x_{9}\,x_{24} -24\,x_{9}\,x_{25} +22\,x_{10} \\*
&\quad{} +14\,x_{10}^{2} -36\,x_{10}\,x_{11} -16\,x_{10}\,x_{25} -15\,x_{11} -13\,x_{11}^{2} -24\,x_{11}\,x_{20} -24\,x_{11}\,x_{25} -8\,x_{20}\,x_{25} \\*
&\quad{} -x_{21} -2\,x_{24} -17\,x_{25} -4\,x_{25}^{2} +x_{26} -x_{44} \\[3pt]
R_{16} &={} 32 +46\,x_{9} -2\,x_{9}^{2} -16\,x_{9}\,x_{10} -8\,x_{9}\,x_{11} -32\,x_{9}\,x_{20} -8\,x_{9}\,x_{24} -16\,x_{9}\,x_{25} +38\,x_{10} \\*
&\quad{} +18\,x_{10}^{2} -72\,x_{10}\,x_{11} -24\,x_{10}\,x_{25} +24\,x_{11} +16\,x_{11}^{2} -48\,x_{11}\,x_{20} -16\,x_{20}\,x_{25} -2\,x_{21} \\*
&\quad{} +8\,x_{24} +2\,x_{26} -2\,x_{30} -4\,x_{30}^{2} +x_{31} -2\,x_{44} \\[3pt]
R_{17} &={} -19\,x_{9} +8\,x_{9}^{2} +8\,x_{9}\,x_{20} +8\,x_{9}\,x_{24} -46\,x_{10} -24\,x_{10}^{2} +24\,x_{10}\,x_{20} +24\,x_{10}\,x_{24} -2\,x_{11} \\*
&\quad{} +4\,x_{20} -4\,x_{20}^{2} -8\,x_{20}\,x_{24} +14\,x_{24} -4\,x_{24}^{2} -x_{25} +x_{32} \\[3pt]
R_{18} &={} -3 -6\,x_{9} -4\,x_{11} -2\,x_{25} -4\,x_{32} -4\,x_{32}^{2} +x_{33} \\[3pt]
R_{19} &={} 8 +42\,x_{9} +7\,x_{9}^{2} +28\,x_{9}\,x_{21} +48\,x_{10} -32\,x_{10}^{2} -32\,x_{10}\,x_{20} +64\,x_{10}\,x_{21} -40\,x_{20}^{2} \\*
&\quad{} +32\,x_{20}\,x_{21} +14\,x_{21} -4\,x_{21}^{2} +x_{34} +2\,x_{37} \\[3pt]
R_{20} &={} 1 +6\,x_{9} -36\,x_{9}^{2} -24\,x_{9}\,x_{34} +2\,x_{34} -4\,x_{34}^{2} +x_{35} \\[3pt]
R_{21} &={} -x_{10}^{2} -4\,x_{10}\,x_{35} -4\,x_{35}^{2} +x_{36} \\[3pt]
R_{22} &={} -6 -2\,x_{9} -4\,x_{10} -14\,x_{20} -16\,x_{20}^{2} +x_{37} \\[3pt]
R_{23} &={} 24 +21\,x_{10} -21\,x_{10}^{2} -18\,x_{10}\,x_{11} -24\,x_{10}\,x_{24} -24\,x_{10}\,x_{37} +3\,x_{11} -9\,x_{11}^{2} -16\,x_{11}\,x_{37} \\*
&\quad{} +3\,x_{24} -12\,x_{24}^{2} -8\,x_{24}\,x_{37} +4\,x_{37} -4\,x_{37}^{2} +x_{38} +2\,x_{44} \\[3pt]
R_{24} &={} 16 -10\,x_{9} -20\,x_{9}^{2} -16\,x_{9}\,x_{11} -24\,x_{9}\,x_{25} -24\,x_{9}\,x_{38} -12\,x_{11} -16\,x_{11}\,x_{25} -16\,x_{11}\,x_{38} \\*
&\quad{} -14\,x_{25} -4\,x_{25}^{2} -8\,x_{25}\,x_{38} +4\,x_{38} -4\,x_{38}^{2} +x_{39} \\[3pt]
R_{25} &={} -x_{11}^{2} -2\,x_{11}\,x_{12} -x_{12}^{2} +x_{40} \\[3pt]
R_{26} &={} -29 +x_{9}^{2} -4\,x_{9}\,x_{40} -14\,x_{11} -2\,x_{11}^{2} -16\,x_{11}\,x_{24} +16\,x_{11}\,x_{25} -8\,x_{11}\,x_{40} -8\,x_{24} \\*
&\quad{} +12\,x_{24}^{2} -8\,x_{24}\,x_{25} +22\,x_{25} -4\,x_{25}^{2} +4\,x_{40} -4\,x_{40}^{2} +x_{41} -4\,x_{44} \\[3pt]
R_{27} &={} 4 +2\,x_{9}^{2} -8\,x_{9}\,x_{40} -4\,x_{11}^{2} -16\,x_{11}\,x_{40} -4\,x_{25} -16\,x_{25}^{2} -32\,x_{25}\,x_{40} -16\,x_{25}\,x_{41} \\*
&\quad{} +4\,x_{40} -24\,x_{40}^{2} -16\,x_{40}\,x_{41} -4\,x_{41}^{2} +x_{42} \\[3pt]
R_{28} &={} -16\,x_{9}^{2} +8\,x_{9}\,x_{11} +32\,x_{9}\,x_{25} +16\,x_{9}\,x_{28} +x_{11} -x_{11}^{2} -8\,x_{11}\,x_{25} -4\,x_{11}\,x_{28} \\*
&\quad{} -16\,x_{25}^{2} -16\,x_{25}\,x_{28} -4\,x_{28}^{2} +x_{43} \\[3pt]
R_{29} &={} -1 -8\,x_{11} -4\,x_{11}^{2} +8\,x_{11}\,x_{24} +8\,x_{24} -4\,x_{24}^{2} +x_{44} \\[3pt]
R_{30} &={} -1 -9\,x_{11}^{2} +24\,x_{11}\,x_{24} +12\,x_{11}\,x_{44} -16\,x_{24}^{2} -16\,x_{24}\,x_{44} -4\,x_{44}^{2} +x_{45} \\[3pt]
R_{31} &={} 1 +x_{10} -x_{10}^{2} -4\,x_{10}\,x_{11} -4\,x_{10}\,x_{45} -4\,x_{11}^{2} -8\,x_{11}\,x_{45} -4\,x_{45}^{2} +x_{46} \\[3pt]
R_{32} &={} -1 -4\,x_{46} -4\,x_{46}^{2} +x_{47} \\[3pt]
R_{33} &={} -1 -x_{11}^{2} -8\,x_{11}\,x_{46} -4\,x_{11}\,x_{47} -16\,x_{46}^{2} -16\,x_{46}\,x_{47} -4\,x_{47}^{2} +x_{48} \\[3pt]
R_{34} &={} 2 +x_{10} -x_{10}^{2} -4\,x_{10}\,x_{11} -4\,x_{10}\,x_{48} -4\,x_{11}^{2} -8\,x_{11}\,x_{48} -4\,x_{48}^{2} +x_{49} \\[3pt]
R_{35} &={} -4\,x_{43}^{2} -8\,x_{43}\,x_{49} -4\,x_{49}^{2} +x_{50} \\[3pt]
R_{36} &={} -9 +7\,x_{9} -x_{9}^{2} -8\,x_{9}\,x_{43} -4\,x_{9}\,x_{50} +24\,x_{43} -16\,x_{43}^{2} -16\,x_{43}\,x_{50} +12\,x_{50} -4\,x_{50}^{2} \\*
&\quad{} +x_{51} \\[3pt]
R_{37} &={} -1 -4\,x_{51} -4\,x_{51}^{2} +x_{52} \\[3pt]
R_{38} &={} -x_{11}^{2} -8\,x_{11}\,x_{51} -4\,x_{11}\,x_{52} -16\,x_{51}^{2} -16\,x_{51}\,x_{52} -4\,x_{52}^{2} +x_{53} \\[3pt]
R_{39} &={} -x_{9} -x_{9}^{2} -4\,x_{9}\,x_{11} -4\,x_{9}\,x_{53} -4\,x_{11} -4\,x_{11}^{2} -8\,x_{11}\,x_{53} -4\,x_{53} -4\,x_{53}^{2} +x_{54} \\[3pt]
R_{40} &={} -10 -12\,x_{54} -4\,x_{54}^{2} +x_{55} \\[3pt]
R_{41} &={} 2 +4\,x_{11} +16\,x_{11}\,x_{54} +8\,x_{11}\,x_{55} -4\,x_{24} -16\,x_{24}\,x_{54} -8\,x_{24}\,x_{55} -x_{44} +8\,x_{54} -16\,x_{54}^{2} \\*
&\quad{} -16\,x_{54}\,x_{55} +4\,x_{55} -4\,x_{55}^{2} +x_{56} \\[3pt]
R_{42} &={} -10 +x_{10} -x_{10}^{2} +8\,x_{10}\,x_{11} -8\,x_{10}\,x_{24} -4\,x_{10}\,x_{56} -16\,x_{11}^{2} +32\,x_{11}\,x_{24} +16\,x_{11}\,x_{56} \\*
&\quad{} -16\,x_{24}^{2} -16\,x_{24}\,x_{56} -12\,x_{49} -4\,x_{49}^{2} -4\,x_{56}^{2} +x_{57} \\[3pt]
R_{43} &={} 1 +2\,x_{10} -2\,x_{10}^{2} +16\,x_{10}\,x_{11} -16\,x_{10}\,x_{24} -8\,x_{10}\,x_{56} -32\,x_{11}^{2} +64\,x_{11}\,x_{24} \\*
&\quad{} +32\,x_{11}\,x_{56} -32\,x_{24}^{2} -32\,x_{24}\,x_{56} -16\,x_{49}^{2} -16\,x_{49}\,x_{57} -8\,x_{56}^{2} -4\,x_{57}^{2} +x_{58} \\[3pt]
R_{44} &={} -6\,x_{9} -3\,x_{9}^{2} -18\,x_{9}\,x_{10} -16\,x_{10} -27\,x_{10}^{2} +4\,x_{10}\,x_{11} +3\,x_{11} +3\,x_{11}^{2} -4\,x_{11}\,x_{20} \\*
&\quad{} -19\,x_{20} -8\,x_{20}^{2} +x_{21} +x_{59} -4\,x_{72} \\[3pt]
R_{45} &={} 1 -8\,x_{9} -4\,x_{9}\,x_{11} +8\,x_{9}\,x_{59} -16\,x_{10} -8\,x_{10}\,x_{11} +16\,x_{10}\,x_{59} +2\,x_{11} +4\,x_{11}\,x_{20} -4\,x_{11}\,x_{58} \\*
&\quad{} +8\,x_{20} -8\,x_{20}\,x_{59} -4\,x_{58}^{2} +3\,x_{59} -4\,x_{59}^{2} +x_{60} \\[3pt]
R_{46} &={} -2 -12\,x_{9} -16\,x_{9}\,x_{11} +16\,x_{9}\,x_{60} -24\,x_{10} -32\,x_{10}\,x_{11} +32\,x_{10}\,x_{60} +2\,x_{11} +6\,x_{11}^{2} \\*
&\quad{} +16\,x_{11}\,x_{20} +8\,x_{11}\,x_{58} +12\,x_{20} -16\,x_{20}\,x_{60} +8\,x_{58}^{2} -4\,x_{59} -2\,x_{60} -4\,x_{60}^{2} +x_{61} \\[3pt]
R_{47} &={} -2 -4\,x_{61} -4\,x_{61}^{2} +x_{62} \\[3pt]
R_{48} &={} -x_{3}\,x_{14} +x_{63} \\[3pt]
R_{49} &={} x_{13} -x_{14} -2\,x_{23}\,x_{63} -x_{64} \\[3pt]
R_{50} &={} x_{13} -x_{14} -2\,x_{27}\,x_{63} -x_{65} \\[3pt]
R_{51} &={} -2\,x_{29}\,x_{63} +2\,x_{42}\,x_{63} -2\,x_{63} -x_{66} +x_{71} \\[3pt]
R_{52} &={} x_{13} -x_{14} -2\,x_{31}\,x_{63} -x_{67} \\[3pt]
R_{53} &={} x_{13} -x_{14} -2\,x_{33}\,x_{63} +6\,x_{63} -x_{68} \\[3pt]
R_{54} &={} x_{13} -x_{14} -4\,x_{35}\,x_{63} -2\,x_{36}\,x_{63} -x_{69} \\[3pt]
R_{55} &={} x_{13} -x_{14} -2\,x_{39}\,x_{63} -x_{70} \\[3pt]
R_{56} &={} x_{13} -x_{14} -2\,x_{42}\,x_{63} -4\,x_{63} -x_{71} \\[3pt]
R_{57} &={} -x_{9}\,x_{10} +x_{9}\,x_{11} +x_{10} +x_{10}\,x_{11} -x_{11} -x_{11}^{2} +x_{72} \\[3pt]
R_{58} &={} 1 +2\,x_{12} +x_{12}^{2} -x_{72}^{2} +x_{73} \\[3pt]
R_{59} &={} -4\,x_{61}\,x_{63} -2\,x_{62}\,x_{63} +2\,x_{63} +x_{74} \\[3pt]
R_{60} &={} -x_{13}^{2} +2\,x_{13}\,x_{14} +2\,x_{13}\,x_{74} -x_{14}^{2} -2\,x_{14}\,x_{74} -x_{73}^{2} -x_{74}^{2} +x_{75} \\[3pt]
R_{61} &={} -1 -x_{3}\,x_{10} +x_{76} \\[3pt]
R_{62} &={} -x_{14}\,x_{76} +x_{78}\,x_{79} +x_{80} \\[3pt]
R_{63} &={} -1 +x_{3} -x_{3}\,x_{9} +x_{77} \\[3pt]
R_{64} &={} -1 +2\,x_{3} +2\,x_{3}\,x_{9} +4\,x_{3}\,x_{10} -2\,x_{3}\,x_{20} +x_{78} \\[3pt]
R_{65} &={} -x_{13}\,x_{77} +x_{79} \\[3pt]
R_{66} &={} -6\,x_{3}\,x_{14} +x_{13}\,x_{76} +12\,x_{14} -8\,x_{14}\,x_{76} -4\,x_{14}\,x_{77} -x_{14}\,x_{78} +x_{81} \\[3pt]
R_{67} &={} -x_{13}\,x_{64} +x_{14}\,x_{64} +x_{82} \\[3pt]
R_{68} &={} -x_{65}\,x_{82} +x_{83} \\[3pt]
R_{69} &={} -x_{66}\,x_{83} +x_{84} \\[3pt]
R_{70} &={} -x_{67}\,x_{84} +x_{85} \\[3pt]
R_{71} &={} -x_{68}\,x_{85} +x_{86} \\[3pt]
R_{72} &={} -x_{69}\,x_{86} +x_{87} \\[3pt]
R_{73} &={} -x_{70}\,x_{87} +x_{88} \\[3pt]
R_{74} &={} -x_{71}\,x_{88} +x_{89} \\[3pt]
R_{75} &={} -x_{80}\,x_{89} +x_{90} \\[3pt]
R_{76} &={} -x_{81}\,x_{90} +x_{91} \\[3pt]
R_{77} &={} -x_{91}^{2} +x_{92} \\[3pt]
R_{78} &={} -x_{2} -x_{2}\,x_{17} +x_{17} +x_{93} \\[3pt]
R_{79} &={} -x_{15}\,x_{93} +x_{94} \\[3pt]
R_{80} &={} -2\,x_{1}^{2} +2\,x_{1}\,x_{13} +2\,x_{1}\,x_{14} +2\,x_{1}\,x_{94} +2\,x_{1}\,x_{95} -x_{13}^{2} -2\,x_{13}\,x_{94} -x_{14}^{2} \\*
&\quad{} -2\,x_{14}\,x_{95} +x_{18}\,x_{97} -x_{75}\,x_{92} -x_{94}^{2} -x_{95}^{2} \\[3pt]
R_{81} &={} -x_{2}\,x_{16} +x_{16} -x_{16}\,x_{93} +x_{95} \\[3pt]
R_{82} &={} -x_{9} -x_{10} -4\,x_{20} -2\,x_{23} -2\,x_{27} -2\,x_{29} -2\,x_{31} -2\,x_{33} -4\,x_{35} -2\,x_{36} -2\,x_{39} -2\,x_{42} -4\,x_{61} \\*
&\quad{} -2\,x_{62} +x_{96} \\[3pt]
R_{83} &={} -x_{98}\,x_{101} +x_{101} +x_{103} \\[3pt]
R_{84} &={} -x_{5}\,x_{125} +x_{99}\,x_{224} -x_{102} -x_{103}\,x_{125} \\[3pt]
R_{85} &={} x_{104}\,x_{225} -x_{106}\,x_{128} -x_{107} \\[3pt]
R_{86} &={} x_{108}\,x_{225} -x_{109}\,x_{128} -x_{110} \\[3pt]
R_{87} &={} x_{111}\,x_{225} -x_{112}\,x_{128} -x_{113} \\[3pt]
R_{88} &={} x_{114}\,x_{225} -x_{115}\,x_{128} -x_{116} \\[3pt]
R_{89} &={} x_{117}\,x_{226} -x_{119}\,x_{128} -x_{120} \\[3pt]
R_{90} &={} x_{121}\,x_{226} -x_{122}\,x_{128} -x_{123} \\[3pt]
R_{91} &={} -x_{18}\,x_{125} +x_{128} \\[3pt]
R_{92} &={} x_{124} -x_{126}\,x_{128} -x_{127} \\[3pt]
R_{93} &={} -x_{128}\,x_{130} +x_{129} -x_{131} \\[3pt]
R_{94} &={} -x_{128}\,x_{133} +x_{132}\,x_{224} -x_{134} \\[3pt]
R_{95} &={} -x_{128}\,x_{136} +x_{135}\,x_{225} -x_{137} \\[3pt]
R_{96} &={} x_{138} -2\,x_{138}\,x_{140} +x_{141} \\[3pt]
R_{97} &={} -x_{138}\,x_{141} -x_{141} +x_{142} \\[3pt]
R_{98} &={} x_{141}^{2} -x_{142}^{2} +x_{162} \\[3pt]
R_{99} &={} -x_{139}\,x_{144} -x_{144} +x_{145} \\[3pt]
R_{100} &={} 2\,x_{143} +x_{143}^{2} +x_{144}^{2} -x_{145}^{2} \\[3pt]
R_{101} &={} 1 -2\,x_{0} +x_{0}^{2} -2\,x_{0}\,x_{139} +x_{146} \\[3pt]
R_{102} &={} -1 -x_{0}\,x_{144} +x_{3} +x_{139}\,x_{144} -x_{143} +x_{144} +x_{146}\,x_{147} \\[3pt]
R_{103} &={} 1 -2\,x_{5} +x_{5}^{2} -2\,x_{5}\,x_{139} +x_{148} \\[3pt]
R_{104} &={} -x_{5}\,x_{193} +x_{18} +x_{148}\,x_{149} -x_{192} +x_{194} \\[3pt]
R_{105} &={} -1 +2\,x_{139}\,x_{208} -x_{150} +2\,x_{208} -x_{208}^{2} \\[3pt]
R_{106} &={} -1 +x_{98} +x_{150}\,x_{151} -x_{192} -x_{193}\,x_{208} +x_{194} \\[3pt]
R_{107} &={} -x_{139}\,x_{153} -x_{153} +x_{154} \\[3pt]
R_{108} &={} 2\,x_{152} +x_{152}^{2} +x_{153}^{2} -x_{154}^{2} \\[3pt]
R_{109} &={} 1 -2\,x_{98} +x_{98}^{2} -2\,x_{98}\,x_{139} +x_{155} \\[3pt]
R_{110} &={} -1 -x_{98}\,x_{153} +x_{99} -x_{152} +x_{154} +x_{155}\,x_{156} \\[3pt]
R_{111} &={} -1 -x_{98}\,x_{144} +2\,x_{98}\,x_{157} +x_{100} -2\,x_{139}\,x_{157} -x_{143} +x_{144} +x_{145} +x_{155}\,x_{157} -3\,x_{157} \\[3pt]
R_{112} &={} 1 +x_{139} -x_{139}\,x_{159} -x_{159} +x_{160} \\[3pt]
R_{113} &={} 1 +2\,x_{158} +x_{158}^{2} -2\,x_{159} +x_{159}^{2} -x_{160}^{2} \\[3pt]
R_{114} &={} -1 +x_{98} -x_{98}\,x_{159} +x_{104} +x_{155}\,x_{161} -x_{158} +x_{160} \\[3pt]
R_{115} &={} -2\,x_{139} -x_{139}^{2} +x_{162} \\[3pt]
R_{116} &={} 1 -x_{98} -x_{98}\,x_{143} -x_{98}\,x_{145} +2\,x_{98}\,x_{163} +x_{105} +x_{139}\,x_{145} -2\,x_{139}\,x_{163} +x_{143} \\*
&\quad{} -x_{144}\,x_{162} +2\,x_{145} +x_{155}\,x_{163} -3\,x_{163} \\[3pt]
R_{117} &={} -x_{139}\,x_{165} -x_{165} +x_{166} \\[3pt]
R_{118} &={} 2\,x_{164} +x_{164}^{2} +x_{165}^{2} -x_{166}^{2} \\[3pt]
R_{119} &={} -1 -x_{98}\,x_{165} +x_{108} +x_{155}\,x_{167} -x_{164} +x_{166} \\[3pt]
R_{120} &={} -x_{139}\,x_{169} -x_{169} +x_{170} \\[3pt]
R_{121} &={} 2\,x_{168} +x_{168}^{2} +x_{169}^{2} -x_{170}^{2} \\[3pt]
R_{122} &={} -1 -x_{98}\,x_{169} +x_{111} +x_{155}\,x_{171} -x_{168} +x_{170} \\[3pt]
R_{123} &={} -x_{139}\,x_{173} -x_{173} +x_{174} \\[3pt]
R_{124} &={} 2\,x_{172} +x_{172}^{2} +x_{173}^{2} -x_{174}^{2} \\[3pt]
R_{125} &={} -1 -x_{98}\,x_{173} +x_{114} +x_{155}\,x_{175} -x_{172} +x_{174} \\[3pt]
R_{126} &={} -x_{139}\,x_{177} -x_{177} +x_{178} \\[3pt]
R_{127} &={} 2\,x_{176} +x_{176}^{2} +x_{177}^{2} -x_{178}^{2} \\[3pt]
R_{128} &={} -1 -x_{98}\,x_{177} +x_{117} +x_{155}\,x_{179} -x_{176} +x_{178} \\[3pt]
R_{129} &={} -x_{139}\,x_{181} -x_{181} +x_{182} \\[3pt]
R_{130} &={} 2\,x_{180} +x_{180}^{2} +x_{181}^{2} -x_{182}^{2} \\[3pt]
R_{131} &={} -1 -x_{98}\,x_{181} +2\,x_{98}\,x_{183} +x_{118} -2\,x_{139}\,x_{183} +x_{155}\,x_{183} -x_{180} +x_{181} +x_{182} -3\,x_{183} \\[3pt]
R_{132} &={} -x_{139}\,x_{185} -x_{185} +x_{186} \\[3pt]
R_{133} &={} 2\,x_{184} +x_{184}^{2} +x_{185}^{2} -x_{186}^{2} \\[3pt]
R_{134} &={} -1 -x_{98}\,x_{185} +x_{121} +x_{155}\,x_{187} -x_{184} +x_{186} \\[3pt]
R_{135} &={} -x_{139}\,x_{189} -x_{189} +x_{190} \\[3pt]
R_{136} &={} 2\,x_{188} +x_{188}^{2} +x_{189}^{2} -x_{190}^{2} \\[3pt]
R_{137} &={} -1 -x_{98}\,x_{189} +x_{124} +x_{155}\,x_{191} -x_{188} +x_{190} \\[3pt]
R_{138} &={} 1 -x_{139}\,x_{193} -x_{193} +x_{194} \\[3pt]
R_{139} &={} -1 +2\,x_{192} +x_{192}^{2} +x_{193}^{2} -2\,x_{194} -x_{194}^{2} \\[3pt]
R_{140} &={} -x_{98}\,x_{193} +2\,x_{98}\,x_{195} +x_{125} -2\,x_{139}\,x_{195} +x_{155}\,x_{195} -x_{192} +x_{193} +x_{194} -3\,x_{195} \\[3pt]
R_{141} &={} -x_{139}\,x_{197} -x_{197} +x_{198} \\[3pt]
R_{142} &={} 2\,x_{196} +x_{196}^{2} +x_{197}^{2} -x_{198}^{2} \\[3pt]
R_{143} &={} -1 -x_{98}\,x_{197} +x_{129} +x_{155}\,x_{199} -x_{196} +x_{198} \\[3pt]
R_{144} &={} -x_{139}\,x_{201} -x_{201} +x_{202} \\[3pt]
R_{145} &={} 2\,x_{200} +x_{200}^{2} +x_{201}^{2} -x_{202}^{2} \\[3pt]
R_{146} &={} -1 -x_{98}\,x_{201} +x_{132} +x_{155}\,x_{203} -x_{200} +x_{202} \\[3pt]
R_{147} &={} -x_{139}\,x_{205} -x_{205} +x_{206} \\[3pt]
R_{148} &={} 2\,x_{204} +x_{204}^{2} +x_{205}^{2} -x_{206}^{2} \\[3pt]
R_{149} &={} -1 -x_{98}\,x_{205} +x_{135} +x_{155}\,x_{207} -x_{204} +x_{206} \\[3pt]
R_{150} &={} -4 -2\,x_{0} -4\,x_{0}^{2} -2\,x_{0}\,x_{4} -x_{1}^{2} -x_{2}^{2} -x_{3}^{2} -x_{4}^{2} -x_{5}^{2} +2\,x_{9} +2\,x_{9}\,x_{10} \\*
&\quad{} -x_{11}^{2} -x_{12}^{2} -x_{13}^{2} -x_{14}^{2} -x_{15}^{2} -x_{16}^{2} -x_{17}^{2} -x_{18}^{2} \\*
&\quad{} -x_{20} -x_{96}^{2} +2\,x_{98} -2\,x_{98}^{2} -x_{99}^{2} -x_{100}^{2} -x_{104}^{2} -x_{105}^{2} \\*
&\quad{} -x_{108}^{2} -x_{111}^{2} -x_{114}^{2} -x_{117}^{2} -x_{118}^{2} -x_{121}^{2} -x_{124}^{2} \\*
&\quad{} -x_{125}^{2} -x_{129}^{2} -x_{132}^{2} -x_{135}^{2} +x_{138} -x_{208}^{2} \\[3pt]
R_{151} &={} -x_{5} -x_{5}\,x_{9} -x_{5}\,x_{10} -x_{5}\,x_{11} -x_{5}\,x_{12} -x_{5}\,x_{13} -x_{5}\,x_{14} -x_{5}\,x_{15} -x_{5}\,x_{16} \\*
&\quad{} -x_{5}\,x_{17} +x_{208} \\[3pt]
R_{152} &={} -16 +x_{0} -x_{209}^{2} \\[3pt]
R_{153} &={} -x_{139}\,x_{211} -x_{211} +x_{220} \\[3pt]
R_{154} &={} -1 +x_{210}^{2} +x_{211}^{2} -x_{220}^{2} \\[3pt]
R_{155} &={} -x_{213} -x_{213}\,x_{214} +x_{221} \\[3pt]
R_{156} &={} -1 +x_{212}^{2} +x_{213}^{2} -x_{221}^{2} \\[3pt]
R_{157} &={} -x_{205}^{2} +x_{222} \\[3pt]
R_{158} &={} x_{211} -2\,x_{216}\,x_{222} -x_{222} \\[3pt]
R_{159} &={} -4\,x_{211} -8\,x_{211}\,x_{215} +x_{214} \\[3pt]
R_{160} &={} -x_{139} -x_{210}\,x_{217} +x_{214} \\[3pt]
R_{161} &={} -1 -x_{204} -x_{210}\,x_{218} +x_{212} \\[3pt]
R_{162} &={} -x_{96} -4\,x_{205}\,x_{219} +x_{213} \\[3pt]
R_{163} &={} -2 +2\,x_{0}^{2} +2\,x_{0}\,x_{4} -2\,x_{0}\,x_{144} -2\,x_{0}\,x_{153} +x_{1}^{2} -2\,x_{1}\,x_{193} +x_{2}^{2} -2\,x_{2}\,x_{181} \\*
&\quad{} +x_{4}^{2} -2\,x_{4}\,x_{153} +x_{9}^{2} -2\,x_{9}\,x_{159} +x_{10}^{2} -2\,x_{10}\,x_{165} +x_{11}^{2} -2\,x_{11}\,x_{169} \\*
&\quad{} +x_{12}^{2} -2\,x_{12}\,x_{173} +x_{13}^{2} -2\,x_{13}\,x_{177} +x_{14}^{2} -2\,x_{14}\,x_{185} +x_{15}^{2} \\*
&\quad{} -2\,x_{15}\,x_{189} +x_{16}^{2} -2\,x_{16}\,x_{197} +x_{17}^{2} -2\,x_{17}\,x_{201} -2\,x_{96}\,x_{205} -x_{139}\,x_{223} \\*
&\quad{} +x_{143}^{2} +x_{152}^{2} +x_{158}^{2} +2\,x_{159} +x_{164}^{2} +x_{168}^{2} +x_{172}^{2} \\*
&\quad{} +x_{176}^{2} +x_{180}^{2} +x_{184}^{2} +x_{188}^{2} +x_{192}^{2} +x_{196}^{2} +x_{200}^{2} \\*
&\quad{} +x_{204}^{2} +2\,x_{205}^{2} -x_{227}^{2} -x_{228}^{2} -x_{229} -x_{229}^{2} \\[3pt]
R_{164} &={} -x_{100}\,x_{224} +x_{125} \\[3pt]
R_{165} &={} -x_{105}\,x_{225} +x_{125} \\[3pt]
R_{166} &={} -x_{118}\,x_{226} +x_{125} \\[3pt]
R_{167} &={} -1 +x_{2} -x_{3}\,x_{234} \\[3pt]
R_{168} &={} -2 -x_{96}^{2} -x_{144}^{2} -x_{153}^{2} +2\,x_{159} -x_{159}^{2} -x_{165}^{2} -x_{169}^{2} \\*
&\quad{} -x_{173}^{2} -x_{177}^{2} -x_{181}^{2} -x_{185}^{2} -x_{189}^{2} -x_{193}^{2} -x_{197}^{2} \\*
&\quad{} -x_{201}^{2} +x_{205}^{2} -x_{227}^{2} -x_{228}^{2} -x_{229} -x_{229}^{2} \\[3pt]
R_{169} &={} -1 -4\,x_{102}^{2} -4\,x_{107}^{2} -4\,x_{110}^{2} -4\,x_{113}^{2} -4\,x_{116}^{2} -4\,x_{120}^{2} \\*
&\quad{} -4\,x_{123}^{2} +x_{125}^{2} -4\,x_{127}^{2} -4\,x_{131}^{2} -4\,x_{134}^{2} -4\,x_{137}^{2} -x_{230}^{2} \\*
&\quad{} -x_{231}^{2} -x_{232} -x_{232}^{2} \\[3pt]
R_{170} &={} -2 +x_{98} -4\,x_{233} -4\,x_{233}^{2} \\[3pt]
R_{171} &={} -1 +x_{234} -4\,x_{235} -4\,x_{235}^{2}
\end{align*}
\endgroup


\section{The compact fixed coefficient specification}\label{app:coefficients}
The compact equation retains its native/K215 ancestry. The direct
quartics in Appendix~\ref{app:direct} use the separate proof-calculus route.
The four coefficient payloads in \texttt{compact/coefficients/} have
bit lengths $590{,}069$, $592{,}884$, $590{,}069$, and $590{,}070$
for $\beta,K_0,M_C,M_F$ respectively. The following complete finite
recipe is retained unchanged and its reconstruction is checked exactly.
% Exact finite recipe extracted from the retained indexed_constants.tex.
% In this appendix K denotes the fixed port K_0, and all other symbols are constants.
\begingroup
\providecommand{\fl}[1]{\left\lfloor#1\right\rfloor}
\providecommand{\rg}[2]{\{#1,\ldots,#2\}}
In this appendix $K$ denotes the port $K_0$ in the new equations.
Every symbol introduced here is a fixed constant, independent of the ten
quantified variables. The formulas and 436-character seed are the retained
indexed BFS943 radix67 specification; they reconstruct the exact supplied
payloads, with $\beta=2^{590069}-1$.

\textbf{The four fixed integers.}
Every index, set, sum, product and floor below defines a constant, not
another unknown. Put
\[
\rho=2^{67},\quad B=\rho^{8807},\quad\beta=B-1,\quad
G_m(x)=\sum_{j=0}^{m-1}x^j.
\]
Define
\[
\begin{split}
\mathcal E={}&\rg0{942}\cup\rg{983}{992}\cup\rg{1003}{1011}\\
&\cup\rg{1013}{1031}\cup\{1063,1064,1095,2195,3296\}.
\end{split}
\]
Then
\begin{align*}
M_C={}&B-1-\sum_{e\in\mathcal E\setminus\{1\}}\rho^e-2\rho^{1063},\\
M_F={}&B-1+4+(2^{51}-2)(\rho^{4393}+\rho^{5495})\\
&+\rho^{1013}G_{10}(\rho)+\rho^{1065}G_9(\rho)
  +\rho^{2195}+\rho^{3296}\\
&+(2^{65}-1)\rho^{8792}G_{13}(\rho)+(2^{19}-1)\rho^{8805},\\
K={}&\rho^{30}+\rho^{62}+\rho^{1100}+\rho^{2201}+\rho^{8849}\\
&+(\rho^{4393}+\rho^{5495})\Phi(1/\rho),
\end{align*}
where
\[
\begin{split}
\Phi(x)={}&(1+2^{51})G_{943}(x)
 -2^{10}\sum_{i=0}^{942}v(t_i)x^i\\
&+\sum_{(b,e,m)\in\mathcal C}2^b x^e G_m(2x),\\
\mathcal C={}&\{(10,1023,8),(18,1013,9),(27,1003,8),\\
&(35,983,9),(44,1031,1),(45,1022,1),(46,1011,1),\\
&(47,992,1),(48,1095,1),(49,2195,1),(50,3296,1)\}.
\end{split}
\]
The rational intermediate in K cancels exactly; all four results are positive integers.

\textbf{The 943 indexed fixed tiles.}
Use the actions $u_j,v_j,d_j$ decoded below. Let
\[
\begin{gathered}
h(q,b)=2(q+1)+b,\quad
\mathcal R=\{u_j:d_j=1,\ 0\le j<430\},\\
\mathcal L=\{u_j:d_j=0,\ 0\le j<430\},\\
r_q=6+|\mathcal R\cap\rg0{q-1}|,\quad
l_q=205+|\mathcal L\cap\rg0{q-1}|.
\end{gathered}
\]
An interval ending at $-1$ is empty. The table defines $t_i$ for
$0\le i\le942$; all index ranges are fixed:
\[
\setlength{\arraycolsep}{1.4pt}\begin{array}{c|l}
\text{index }i&\text{tile }t_i=(w,n,e,s)\\\hline
0,1&(2,435,3,1),\ (5,0,0,0)\\
j+2\ (j=2q+b<430)&((1-d_j)l_{u_j},h(q,b),d_jr_{u_j},v_j)\\
432+b\ (b=0,1)&(0,h(215,b),0,h(215,b))\\
434+b\ (b=0,1)&(0,b,0,b)\\
424+2r_q+b\ (q\in\mathcal R,b=0,1)&(r_q,b,0,h(q,b))\\
424+2l_q+b\ (q\in\mathcal L,b=0,1)&(0,b,l_q,h(q,b))\\
924+j\ (0\le j<4)&(0,j+430\fl{j/2},0,438)\\
928&(2,434,2,0)\\
929&(3,436,4,4)\\
930&(4,436,4,0)\\
497+t\ (t=434,435,436)&(1,438,1,t)\\
\substack{934+3a+b-6\fl{a/4}\\a\in\{0,1,4\},\ b=0,1,2}&(a,437,b,437)
\end{array}
\]
For $t=(w,n,e,s)$ put
\[
v(t)=w+2^8n+2^{17}e+2^{25}s+127\,2^{34}.
\]

\textbf{A closed decoder for the 430 actions.}
Put $D_b(A,k)=\fl{A/b^k}\bmod b$, where
$a\bmod b=a-b\fl{a/b}$. For $0\le q<215$, $b=0,1$, $j=2q+b$:
\[
\begin{gathered}
N=\mathcal A\bmod2^{636},\quad
T_0=\fl{\mathcal A/2^{636}}\bmod2^{1675},\\
E_0=\fl{\mathcal A/2^{2311}},\\
\nu_j=D_2(N,j),\quad \sigma_j=\sum_{k=0}^{j-1}\nu_k,\\
m_q=D_2(N,430+q),\quad \tau_q=\sum_{k=0}^{q-1}m_k,\\
F_q=14(1-m_q)+m_qD_{16}(E_0,\tau_q),\\
u_j=\nu_j(1+\sigma_j)+(1-\nu_j)D_{216}(T_0,j-\sigma_j),\\
v_j=D_2(F_q,2b),\quad d_j=D_2(F_q,2b+1).
\end{gathered}
\]
Empty sums are zero. Here $(v_0,d_0,u_0)=(1,1,1)$; the constructor tile
uses this actual first target, including its state renaming.

\textbf{The integer $\mathcal A$.}
The alphabet, in order of digit values $0,\ldots,61$, is
\par\noindent
{\fontsize{9.0}{10.5}\selectfont\ttfamily 0123456789ABCDEFGHIJKLMNOPQRSTUVWXYZabcdefghijklmnopqrstuvwxyz}.\par
Concatenate the following 436 characters as digits $e_0,\ldots,e_{435}$;
put $\mathcal A=\sum_{j=0}^{435}e_j62^{435-j}$.\par
\begingroup\fontsize{9.0}{10.5}\selectfont\ttfamily
BxEnhvy6NipzkWmROo0GsA3uK3twWLxscLq20UTumy5Cdu461LwD6g7RDNzF8Y\par
iYDFTpB8ehAypd7d6OsOL4APamnJeNXQypUezCmW9n4oXP6XfgmiWuAWZ8ZJ7U\par
OP7gu4kVdlGZqZHULQMar1dXNeJW4i7oNVI8usp8BlovjjZbwPCovZI328fRUq\par
4q92R81s5GjWQaQfGhkq9UYMK3rRMaBmsOIBkJqc2cj2VRIZ6mjzkBIOkW9SjZ\par
FzeA6J5u4j3sBL3SdgwXzM9mHvQivu8EJ1TBTpRhHDqZIbmYgwQ8VqTAFotl6F\par
We1YC554GKrVHsA8oUGKSA6F6k5KREBSD0Tgd7bFALY84LLbpw574BQyvdoHrl\par
0Oe35uwDf15DpT6VJWZBjSbHPpPAINYTte0eUxnTFTKBlGj8Am6sb1gLldh8WZ\par
8P\par
\endgroup

\endgroup



\section{Exact compact degree certificates}
Set $q_0=\beta J$ and $k_0=\eta+\zeta$. The complete highest homogeneous part
of the principal equation is
\[
[P_{\mathrm{tr}}]_{232}
=-9\,2^{10}q_0^{100}w^{26}s^{58}i^{10}k_0^{32}
\bigl(k_0sT-q_0^2(q_0-F)\bigr)^2.
\]
It is nonzero because $\beta>0$. The new degrees of $f,V,Q,G,M,Z$
are respectively $36,40,80,50,11,57$. Thus $-QM^2GV^2$ has degree
$80+22+50+80=232$, while $QZ^2$ has degree $194$ and the remaining
terms are lower. The verifier computes the full highest forms, including
every possible cancellation in the arithmetic circuit.

For each circuit, the following is one coefficient-certified term of
its highest homogeneous part. Here $\beta$ is a fixed positive integer
and is assigned degree zero.
\begin{align*}
d=402:\quad & 1024\,J^{244}\,T^{4}\,\beta^{244}\,\eta^{40}\,i^{12}\,s^{70}\,w^{32}\\
d=352:\quad & -256\,J^{216}\,T^{2}\,\beta^{216}\,\eta^{34}\,i^{10}\,s^{62}\,w^{28}\\
d=290:\quad & -256\,J^{158}\,T^{2}\,\beta^{158}\,\eta^{34}\,i^{10}\,s^{60}\,w^{26}\\
d=258:\quad & -256\,J^{132}\,\beta^{132}\,\eta^{32}\,i^{10}\,s^{58}\,w^{26}\\
d=232:\quad & -9216\,J^{106}\,\beta^{106}\,\eta^{32}\,i^{10}\,s^{58}\,w^{26}
\end{align*}


\begin{thebibliography}{9}
\bibitem{repo} ProveIt contributors, \emph{DiophantineConsistency research
collection}, commit
\texttt{b7e35224c04fe1031c0a53f80a395e9b32205022}, 25 October 2026.
\href{https://github.com/VladimirReshetnikov/ProveIt/tree/b7e35224c04fe1031c0a53f80a395e9b32205022/Computability/DiophantineConsistency}{Pinned research collection on GitHub}.
\bibitem{lucas} AI-assisted ProveIt research, \emph{Lucas corrections and
witness elimination}, incoming archive dated 25 October 2026; exact bytes,
proof, and coefficient closure retained under \texttt{compact/}.
\bibitem{factorialfree} AI-assisted ProveIt research, \emph{Removing factorial
authenticity: binomial prime-power certificates for Diophantine consistency},
incoming archive dated 25 October 2026; source and proof under \texttt{direct/}.
\bibitem{sharederror} AI-assisted ProveIt research, \emph{Shared factorial
error}, incoming archive dated 25 October 2026; source and proof under
\texttt{direct/}.
\bibitem{farhi} B. Farhi, \emph{An identity involving the least common multiple
of binomial coefficients and its application}, American Mathematical
Monthly \textbf{116} (2009), 836--839.
\url{https://arxiv.org/abs/0906.2295}.
\end{thebibliography}
\end{document}
